\begin{recipeDP}
	[
		preparationtime = {\SI{45}{\minute}},
		bakingtime = {\SI{30}{\minute}},
		bakingtemperature = {\protect\bakingtemperature{fanoven=\SI{180}{\celsius}}},
		portion = {6 Portionen},
		source = {@zuckerjagdwurst}
	]
	{\addtoidx{Polenta}-\addtoidx{Lasagne}}

	\graph
		{
			big=Hauptgerichte/Ofengerichte/Polenta-Lasagne/big.jpg,
			small=Hauptgerichte/Ofengerichte/Polenta-Lasagne/small.jpg
		}

	\ingredients{
		\multicolumn{2}{l}{\textbf{Polenta}} \\
		\SI{500}{\g} & Maisgrieß \\
		\SI{2}{\EL} & Hefeflocken \\
		\SI{1050}{\ml} & Wasser \\
		 & Salz \\
		\\
		\multicolumn{2}{l}{\textbf{Tomaten-Hack-Füllung}} \\
		\SI{150}{\g} & Soja-Granulat \\
		\SI{375}{\ml} & Wasser \\
		1 & Zwiebel \\
		2 & Knoblauchzehen \\
		\SI{400}{\g} & Karotten \\
		\SI{4}{\EL} & Tomatenmark \\
		\SI{4}{\EL} & Sojasoße \\
		\SI{2}{\TL} & Paprika edelsüß \\
		\SI{1}{\EL} & Agavendicksaft \\
		\SI{100}{\ml} & Rotwein (oder Balsamico) \\
		\SI{800}{\g} & passierte Tomaten \\
		\SI[parse-numbers = false]{\nicefrac{1}{2}}{\TL} & Oregano \\
		\SI[parse-numbers = false]{\nicefrac{1}{2}}{\TL} & Rosmarin \\
		\SI[parse-numbers = false]{\nicefrac{1}{2}}{\TL} & Thymian \\
		 & Pflanzenöl \\
		\\
		\multicolumn{2}{l}{\textbf{Bohnen-Béchamel}} \\
		\SI{400}{\g} & weiße Bohnen \\
		\SI{4}{\EL} & Butter \\
		\SI{400}{\ml} & Milch \\
		\SI{3}{\EL} & Weizenmehl \\
		\SI{4}{\EL} & Hefeflocken \\
		 & Salz \\
		 & Pfeffer \\
		 & Muskat \\
		\\
        12 & Lasagneplatten \\
		\SI{150}{\g} & Streukäse
	}

	\preparation{
		\step \SI{1050}{\ml} leicht gesalzenes Wasser aufkochen, Maisgrieß und \SI{2}{\EL} Hefeflocken einrühren und 2 Minuten köcheln lassen.
        In eine geölte Auflaufform (dieselbe Form, in der die Lasagne gebacken wird) geben und 2 bis 3 Stunden stehen lassen, bis die Masse gestockt ist.
		\step Zwiebel, Knoblauch und Karotten fein würfeln und 5 Minuten in etwas Öl anbraten, bis die Zwiebeln glasig sind.
        Granulat hinzugeben, bis es gebräunt ist, dann mit \SI{4}{\EL} Sojasoße und \SI{2}{\TL} Paprikapulver würzen.
        Knoblauch zugeben und 2 Minuten mitbraten.
        Tomatenmark kurz anrösten, mit Rotwein und Agavendicksaft ablöschen.
        Sobald der Wein etwas verkocht ist, passierte Tomaten, Ketchup und Gewürze zugeben.
        Mindestens 10 Minuten köcheln lassen und mit Salz und Pfeffer abschmecken.
		\step Vegane Butter schmelzen, Mehl einrühren und 1 Minute bei kleiner Hitze anschwitzen.
        Die Milch langsam unter ständigem Rühren dazugeben. 
        Die Bohnen hinzugeben und glatt pürieren.
        Mit Hefeflocken, Salz, Pfeffer und Muskat abschmecken.
        Bei Bedarf mit mehr pflanzlicher Milch verdünnen.
		\step Gestockte Polenta aus der Form stürzen und je nach Dicke in 3 Schichten teilen.
        In der Auflaufform abwechselnd Bolognese, Polentaplatten und Béchamel schichten und wiederholen, bis Zutaten oder Form voll sind. Bei großer Form mit Lasagneplatten strecken.
        Als oberste Schicht Polenta mit veganem Reibekäse toppen.
		Bei \SI{180}{\celsius} Umluft etwa 30 Minuten backen, bis der Käse goldbraun ist.
	}

	% \suggestion[TITEL EINES VORSCHLAGS]{
	% 	VORSCHLAG (DURCH HORIZONTALE LINIE VOM REZEPT GETRENNT)
	% }
	%
	% \hint{
	%     HINWEIS (IN EINEM KASTEN UNTEN AUF DER SEITE)
	% }

\end{recipeDP}
